%hw-1.tex
%%%%%%%%%%%%%%%%%%%%%%%%%%%%%%%%%%%%%%%%%%%%%%%%%%%%%%%%%%%%%%%%%%%%%%
%%%%%%            Math 403 Homework \#1, Spring 2025            %%%%%%
%%%%%%                                                          %%%%%%
%%%%%%                 Instructor: Ezra Miller                  %%%%%%
%%%%%%%%%%%%%%%%%%%%%%%%%%%%%%%%%%%%%%%%%%%%%%%%%%%%%%%%%%%%%%%%%%%%%%
\documentclass[11pt]{article}

\oddsidemargin=17pt \evensidemargin=17pt
\headheight=9pt     \topmargin=26pt
\textheight=564pt   \textwidth=433.8pt
%voffset=-6ex

\usepackage{amsmath,amssymb,graphicx,color,enumitem}%,setspace,mathrsfs}%,hyperref
\setlist[enumerate]{parsep=1.5ex plus 4pt,topsep=-1ex plus 4pt,itemsep=0ex}
  \setlist[itemize]{parsep=1.5ex plus 4pt,topsep=-1ex plus 4pt,itemsep=-1.33ex}

\leftmargini=5.5ex
\leftmarginii=3.5ex

%for \marginpar to fit optimally
%hoffset=-1.02in
\setlength\marginparwidth{2.2in}
\setlength\marginparsep{1mm}
\newcommand\red[1]{\marginpar{\linespread{.85}\sf%
		   \vspace{-1.4ex}\footnotesize{\color{red}#1}}}
\newcommand\score[1]{\marginpar{\vspace{-2ex}\color{blue}{#1/3}}\hspace{-1ex}}
\newcommand\extra[1]{\marginpar{\color{blue}{#1/1}}\hspace{-1ex}}
\newcommand\total[2]{\marginpar{\colorbox{yellow}{\huge #1/#2}}}
\newcommand\collab[1]{\marginpar{\vspace{-11ex}\colorbox{yellow}{#1/3}}\hspace{-1ex}}
\newcommand\magenta[1]{\colorbox{magenta}{\!\!#1\!\!}}
\newcommand\yellow[1]{\colorbox{yellow}{\!\!#1\!\!}}
\newcommand\green[1]{\colorbox{green}{\!\!#1\!\!}}
\newcommand\cyan[1]{\colorbox{cyan}{\!\!#1\!\!}}
\newcommand\rmagenta[2][0ex]{\red{\vspace{#1}\magenta{\phantom{:}}\,: #2}}
\newcommand\ryellow[2][0ex]{\red{\vspace{#1}\yellow{\phantom{:}}\,: #2}}
\newcommand\rgreen[2][0ex]{\red{\vspace{#1}\green{\phantom{:}}\,: #2}}
\newcommand\rcyan[2][0ex]{\red{\vspace{#1}\cyan{\phantom{:}}\,: #2}}

%For separated lists with consecutive numbering
\newcounter{separated}

\newcommand{\excise}[1]{}
\newcommand{\comment}[1]{{$\star$\sf\textbf{#1}$\star$}}

%%%%%%%%%%%%%%%%%%%%%%%%%%%%%%%%%%%%%%%%%%%%%%%%%%%
%                                                 %
% TO ENABLE GRADING, DO NOT ALTER ABOVE THIS LINE %
%                                                 %
%%%%%%%%%%%%%%%%%%%%%%%%%%%%%%%%%%%%%%%%%%%%%%%%%%%

%new math symbols taking no arguments
\newcommand\0{\mathbf{0}}
\newcommand\CC{\mathbb{C}}
\newcommand\NN{\mathbb{N}}
\newcommand\QQ{\mathbb{Q}}
\newcommand\RR{\mathbb{R}}
\newcommand\ZZ{\mathbb{Z}}
\newcommand\bb{\mathbf{b}}
\newcommand\kk{\Bbbk}
\newcommand\mm{\mathfrak{m}}
\newcommand\xx{\mathbf{x}}
\newcommand\yy{\mathbf{y}}
\newcommand\minus{\smallsetminus}
\newcommand\goesto{\rightsquigarrow}

%redefined math symbols taking no arguments
\newcommand\<{\langle}
\renewcommand\>{\rangle}
\renewcommand\iff{\Leftrightarrow}
\renewcommand\implies{\Rightarrow}

%to explain about LaTeX commands:
\newcommand\command[1]{\texttt{$\backslash$#1}}

%new math symbols taking arguments
\newcommand\ol[1]{{\overline{#1}}}

%redefined math symbols taking arguments
\renewcommand\mod[1]{\ (\mathrm{mod}\ #1)}

%roman font math operators
\DeclareMathOperator\im{im}

%for easy 2 x 2 matrices
\newcommand\twobytwo[1]{\left[\begin{array}{@{}cc@{}}#1\end{array}\right]}

%for easy column vectors of size 2
\newcommand\tworow[1]{\left[\begin{array}{@{}c@{}}#1\end{array}\right]}


%%%%%%%%%%%%%%%%%%%%%%%%%%%%%%%%%%%%%%%%%%%%%%%%%%%%%%%%%%%%%%%%%%%%%%
\begin{document}%%%%%%%%%%%%%%%%%%%%%%%%%%%%%%%%%%%%%%%%%%%%%%%%%%%%%%
%%%%%%%%%%%%%%%%%%%%%%%%%%%%%%%%%%%%%%%%%%%%%%%%%%%%%%%%%%%%%%%%%%%%%%


\title{\mbox{}\\[-8ex]Math 403 Homework \#1, Spring 2025\\\normalsize
	Instructor: Ezra Miller\\[-2.5ex]}
\author{Solutions by: ...your name...\\[1ex]
Collaborators: ...list those with whom you worked on this assignment...\\
(1 point for each of up to 3 collaborators who also list you)}
\date{Due: 11:59pm Saturday 25 January 2025}

\maketitle

\vspace{-5ex}\collab{}%

\noindent
\textsc{Reading assignments}\total{}{57}% cover Lectures 1 - 7
\hfill
(item numbers are lecture numbers as in PDF lecture notes)
\begin{enumerate}[itemsep=-1ex]
\item% 1.
for Thu.~9 January
  \begin{itemize}
  \item{}%
  [Hefferon, p.153--154] field axioms
  \item{}%
  [Cornell, ``Fields'' (about 1/3 of the way through
  \texttt{4330-week1.pdf})]: more on fields%
  \item{}%
  [Climenhaga, \S 7.2, \S 9.1] isomorphism, rank-nullity
  \end{itemize}

\item% 2.
for Tue.~14 January
  \begin{itemize}
  \item{}%
  [Lax, Chapter~1] quotients
  \item{}%
  [Climenhaga, \S 5.1] quotients
  \item{}%
  [Cornell, ``Quotient Spaces'' (\texttt{4330-week5.pdf})] universal properties
  \item{}%
  [Cornell, ``Exact Sequences'' (about halfway through \texttt{4330-week5.pdf})]
  \end{itemize}

\item% 3.
for Thu.~16 January
  \begin{itemize}
  \item{}%
  [Lax, Chapter~2] duality
  \item{}%
  [Climenhaga, \S 5.1, \S 8.2] duality, transpose
  \item{}%
  [Treil, \S 5.1--5.5] Hermitian inner products, ajoints; \S
  5.2--5.4 should be mostly review
  \end{itemize}

\item% 4.
for Tue.~21 January
  \begin{itemize}
  \item{}%
  [Lax, Appendix 15] Jordan canonical form (very short proof)
  \item{}%
  [Hefferon, \S 5.IV.1] characteristic polynomial and minimal polynomial
  \item{}%
  [Treil, \S 9.3 and \S 9.5] generalized eigenspaces and Jordan canonical form
  \end{itemize}

\item% 5.
for Thu.~23 January
  \begin{itemize}
  \item{}%
  [Lax, Chapter 14, p.214--221] norms, equivalence, continuity, local compactness
  \item{}%
  [Stewart--Sun, \S II.1] norms, equivalence, Hahn--Banach theorem
  \end{itemize}

\item% 6.
for Tue.~28 January
  \begin{itemize}
  \item{}%
  [Lax, Chap.12, p.187--190] Convex sets
  \end{itemize}

\item% 7.
for Thu.~30 January
  \begin{itemize}
  \item{}%
  [Serge Lang, \emph{Linear Algebra}, Chapter XII, \S1--\S2]
  separating hyperplanes
  \item{}%
  [Serge Lang, \emph{Linear Algebra}, Chapter XII, \S3--\S4]
  support hyperplane, extreme~point$\!\!$
  \end{itemize}
\end{enumerate}
  
\pagebreak

\noindent
\textsc{Exercises}% cover Lectures 1 - 4

\begin{enumerate}
\item\ \score{}%1.
(Freshman's Dream): The \emph{characteristic} of a field~$F$ is the
smallest positive integer $p$ such that the sum $1 + \cdots + 1$ of
$p$ multiplicative identities is $0$ in~$F$.  Prove that $p$ is prime
if it is finite.  If $F$ has characteristic~$p$, show that $(a + b)^p
= a^p + b^p$ for $a,b \in F$.

\item\ \score{}%2.
Fix a vector space $V$ and a subspace $W \subseteq V$ over a
field~$F$.  Let $\pi: V \to V/W$ be the \emph{projection homomorphism}
given by $\pi(v) = v + W$.  Write $X$ for the set of all subspaces
of~$V$ that contain~$W$, and write $Y$ for the set of all subspaces
of~$V/W$.  Prove that $\pi$ induces a bijection between these two
sets, with
\begin{align*}
X &\to Y
\\
L &\mapsto \pi(L) = \{\pi(v) \mid v \in L\}\\
\intertext{and}\\[-8.75ex]
Y &\to X
\\
M &\mapsto \pi^{-1}(M) = \{v \in V \mid \pi(v) \in M\}.
\end{align*}

\item\ \score{}%3.
Show that giving an exact sequence $\cdots \to V_{i-1} \to V_i \to
V_{i+1} \to \cdots$ is the same as giving a collection of short exact
sequences $0 \to K_i \to V_i \to K_{i+1} \to 0$, one for each~$i$.
(The long exact sequence is said to be constructed by \emph{splicing}
the short exact sequences together.)

\item\ \score{}\label{rn}%4.
Rank-nullity theorem for exact sequences: Given an exact sequence
$$%$$ $
  0 \to V_0 \to V_1 \to \cdots \to V_r \to 0,
$$%$$ $
prove that $\sum_{i=0}^r (-1)^i \dim V_i = 0$.

\item\ \score{}%5.
Rank-nullity for arbitrary complexes: Given a complex $0 \to V_0 \to
V_1 \to \cdots \to V_r \to 0$, write $B_i = \im(V_{i-1} \to V_i)$ and
$Z_i = \ker(V_i \to V_{i+1})$.  Prove that
$$%$$ $
  \sum_{i=0}^r (-1)^i \dim V_i = \sum_i (-1)^i \dim H_i,
$$%$$ $
where $H_i = Z_i/B_i$ is the $i^\text{th}$ homology of the complex.
[In Exercise~\ref{rn}, $B_i = Z_i = K_i$, so $H_i = Z_i/B_i = 0$ for
all~$i$.]

\item\ \score{}%6.
Two elements $u$ and~$v$ in a vector space~$V$ are \emph{congruent
modulo} a subspace~$W \subseteq V$, written $u \equiv v \mod W$, if $u
+ W = v + W$.  Show that congruence modulo~$W$ is an \emph{equivalence
relation} on~$V$, meaning that it is
\begin{itemize}
\item
reflexive: $v \equiv v \mod W$ for all $v \in V$;
\item
symmetric: if $u \equiv v \mod W$, then $v \equiv u \mod W$; and
\item
transitive: if $u \equiv v \mod W$ and $v \equiv x \mod W$ then $u
\equiv x \mod W$.
\end{itemize}

\item\ \score{}%7.
Give an example of three subspaces $Y_1$, $Y_2$, and $Y_3$ in $\RR^2$
such that $Y_1 + Y_2 + Y_3 = \RR^2$ and $Y_i \cap Y_j = \{\0\}$ for
all $i \neq j$, but $\RR^2$ is not the direct sum of $Y_1$, $Y_2$,
and~$Y_3$.

\item\ \score{}%8.
Prove that if $V$ is a vector space and $W \subseteq V$ is a subspace,
then $W$ has a \emph{complement}: a subspace $U \subseteq V$ such that
$V = W \oplus U$.  Hint: $V/W$ has a basis; lift it back to~$V$.

\item%9.
For vectors $\xx = (1, 2i, 1+i)$ and $\yy = (i, 2-i, 3)$, compute
\begin{enumerate}[itemsep=-1ex]
\item\ \score{}%
$\<\xx,\yy\>$, $||\xx||^2$, $||\yy||^2$, and $||\yy||$;
\item\ \score{}%
$\<3\xx,2i\yy\>$ and $\<2\xx,i\xx + 2\yy\>$;
\item\ \score{}%
$||\xx+2\yy||$.  [Use parts (a) and~(b) for this.]
\end{enumerate}

\item\ \score{}%10.
Prove that for vectors in an inner product space,
$||\xx\pm\yy||^2 = ||\xx||^2 + ||\yy||^2 \pm 2\mathrm{Re}\<\xx,\yy\>$.

\item\ \score{}%11.
For any $m \times n$ complex matrix~$A$, prove that $\ker(A^* A) =
\ker(A)$.

\item\ \score{}%12.
Prove that if $P$ is self-ajoint (that is, $P^* = P$) and idempotent
(that is, $P^2 = P$) then $P$ is the matrix for an orthogonal
projection.

\item\ \score{}%13.
If $V$ is a vector space over~$\CC$ of dimension~$n$, then it is also
a vector space over~$\RR$ of dimension~$2n$.  (If this isn't clear to
you, then write down a proof.)  Given a Hermitian inner product
$\<\xx,\yy\>$ on~$V$ as a complex vector space, show that the real
part $\<\xx,\yy\>_\RR = \mathrm{Re}(\<\xx,\yy\>)$ is an inner product
on~$V$ as a real vector space.

\item\ \score{}%14.
Find all possible Jordan forms of linear transformations with
characteristic polynomial $(t - 1)^2 (t + 2)^2$.

\item\ \score{}%15.
Find all possible Jordan forms of linear transformations with
characteristic polynomial $(t - 2)^3 (t + 1)$ and minimal polynomial
$(t - 2)^2 (t + 1)$.

\item\ \score{}%16.
How many similarity classes are there for $3 \times 3$ matrices whose
only eigenvalues are $-3$ and~$4$?

\item\ \score{}%17.
Prove or disprove: two $n \times n$ matrices are similar if and only
if they have the same characteristic polynomial and minimal
polynomial.
\end{enumerate}


%%%%%%%%%%%%%%%%%%%%%%%%%%%%%%%%%%%%%%%%%%%%%%%%%%%%%%%%%%%%%%%%%%%%%%
\end{document}%%%%%%%%%%%%%%%%%%%%%%%%%%%%%%%%%%%%%%%%%%%%%%%%%%%%%%%%
%%%%%%%%%%%%%%%%%%%%%%%%%%%%%%%%%%%%%%%%%%%%%%%%%%%%%%%%%%%%%%%%%%%%%%


%%% Various Emacs customizations:
%%% Local Variables:
%%% fill-column: 70
%%% indent-tabs-mode: t
%%% TeX-electric-sub-and-superscript: nil
%%% TeX-brace-indent-level: 0
%%% LaTeX-indent-level: 0
%%% LaTeX-item-indent: 0
%%% TeX-master: t
%%% End: